\documentclass[aspectratio=169]{beamer}
\input{header}
\coursecode{JKSSB}

\title{File Formats \& Named Technology Associations}
\subtitle{Lesson 55 --- Format and Product Matching}
\author{JKSSB Computer Awareness}
\date{\today}

\begin{document}
\frame{\titlepage}

\begin{frame}{Why Formats Matter}
  \begin{itemize}
    \item Every file on a computer is stored in some \key{format} --- a fixed
      way of arranging its data.
    \item The format decides which program can open, read, and edit the file.
    \item An \key{image format} stores a picture; a \key{document format}
      stores a printable document.
    \item Matching an extension to the right file type and program is a
      high-yield exam skill.
  \end{itemize}
  \begin{block}{The one idea}
    Match \key{extension} $\rightarrow$ \key{format} $\rightarrow$
    \key{software} --- and do not confuse the three.
  \end{block}
\end{frame}

\begin{frame}{Extension and Format Are Not the Same}
  \begin{itemize}
    \item The file name ends with an \key{extension}: the letters after the
      last dot.
    \item Example: in \texttt{report.png}, \key{png} is the extension.
    \item The \key{format} is the way the data inside the file is actually
      stored.
    \item The extension usually signals the format, but the two words mean
      different things.
  \end{itemize}
  \begin{alertblock}{Trap alert}
    ``Extension = format'' is wrong. \texttt{.png} is the \key{extension};
    \key{PNG} is the \key{format} (TRM-09, TRAP-34).
  \end{alertblock}
\end{frame}

\begin{frame}{Image Formats}
  \begin{itemize}
    \item \key{PNG} = Portable Network Graphics, an \key{image} (raster)
      format.
    \item PNG uses \key{lossless} compression and supports transparency.
    \item Other common image formats: \key{JPEG} (\texttt{.jpg/.jpeg}, lossy)
      and \key{GIF} (\texttt{.gif}).
    \item Image extensions are \texttt{.png}, \texttt{.jpg}, \texttt{.jpeg},
      \texttt{.gif}, \texttt{.bmp}.
  \end{itemize}
  \begin{alertblock}{Trap alert}
    PNG is an \key{image format}, never a document, spreadsheet, or database
    format.
  \end{alertblock}
\end{frame}

\begin{frame}{PDF: Portable Document Format}
  \begin{itemize}
    \item \key{PDF} = \key{Portable Document Format}, created by Adobe.
    \item It keeps the exact \key{layout, fonts, and formatting} on every
      device.
    \item Because the page is fixed, it is called a \key{fixed-layout}
      format.
    \item A PDF may contain real text or be a \key{scanned image} of a page.
  \end{itemize}
  \begin{block}{Why PDF is used}
    PDF is the standard for sharing a document that must look the same
    everywhere it is opened.
  \end{block}
\end{frame}

\begin{frame}{PDF Behaviour and Readers}
  \begin{itemize}
    \item A PDF opens in a \key{PDF reader}, such as \key{Adobe Acrobat
      Reader}.
    \item A \key{web browser} can also display a PDF.
    \item A \key{scanned} PDF is a picture; it needs \key{OCR} to make its
      text searchable and selectable.
    \item PDF/A is an archival PDF standard for long-term preservation.
  \end{itemize}
  \begin{alertblock}{Trap alert}
    A PDF is \key{not} an editable Word document. Reading and editing are
    different jobs.
  \end{alertblock}
\end{frame}

\begin{frame}{The Office Document Family}
  \begin{center}
    \begin{tabular}{p{0.16\textwidth}p{0.34\textwidth}p{0.36\textwidth}}
      \toprule
      \textbf{Extension} & \textbf{Format} & \textbf{Program} \\
      \midrule
      \texttt{.docx} & Word document & Microsoft Word \\
      \texttt{.xlsx} & Excel workbook & Microsoft Excel \\
      \texttt{.pptx} & PowerPoint presentation & Microsoft PowerPoint \\
      \bottomrule
    \end{tabular}
  \end{center}
  \vspace{0.25cm}
  \muted{These are the modern Office Open XML formats; the older extensions
  were \texttt{.doc}, \texttt{.xls}, and \texttt{.ppt}.}
\end{frame}

\begin{frame}{Plain-Data Formats: CSV and TXT}
  \begin{itemize}
    \item \key{CSV} = \key{Comma-Separated Values}.
    \item A CSV file stores a table as \key{plain text}, with values
      separated by commas.
    \item \key{TXT} is a plain text file with \key{no formatting}.
    \item Both open in a simple text editor and are not tied to one company.
  \end{itemize}
  \begin{block}{Keep them apart}
    CSV looks like a table but is only text; TXT holds characters with no
    table structure at all.
  \end{block}
\end{frame}

\begin{frame}{Microsoft Access: \texttt{.mdb} vs \texttt{.accdb}}
  \begin{itemize}
    \item Both \texttt{.mdb} and \texttt{.accdb} are file extensions for a
      \key{Microsoft Access database}.
    \item \texttt{.mdb} is the \key{older} Access database format (Access 2003
      and earlier).
    \item \texttt{.accdb} is the \key{newer} Access database format (Access
      2007 and later).
    \item Both store database objects such as tables, queries, forms, and
      reports.
  \end{itemize}
  \begin{alertblock}{Trap alert}
    \texttt{.mdb} and \texttt{.accdb} are \key{both Access databases}, not
    Word or Excel files.
  \end{alertblock}
\end{frame}

\begin{frame}{Named Technology: Mozilla and Firefox}
  \begin{itemize}
    \item \key{Firefox} is a \key{web browser} developed by \key{Mozilla}.
    \item \key{Mozilla} is the organisation; \key{Firefox} is its browser
      product.
    \item Mozilla also makes \key{Thunderbird}, an email client.
    \item \muted{Other browsers: Google Chrome, Microsoft Edge, Safari,
      Opera.}
  \end{itemize}
  \begin{alertblock}{Trap alert}
    Firefox is by \key{Mozilla}, not Microsoft. Microsoft's browser is
    \key{Edge}.
  \end{alertblock}
\end{frame}

\begin{frame}{Named Technology: C-DAC and PARAM}
  \begin{itemize}
    \item \key{C-DAC} = \key{Centre for Development of Advanced Computing}.
    \item C-DAC is an R\&D organisation under \key{MeitY}, Government of
      India.
    \item \key{PARAM} is a family of \key{supercomputers} developed by C-DAC.
    \item \key{PARAM 8000} was India's first indigenous supercomputer.
  \end{itemize}
  \begin{block}{The pair to remember}
    \key{PARAM} is the Indian supercomputer; \key{C-DAC} is the organisation
    that developed it.
  \end{block}
\end{frame}

\begin{frame}{Software Roles: Match the Tool to the Job}
  \begin{center}
    \begin{tabular}{p{0.30\textwidth}p{0.56\textwidth}}
      \toprule
      \textbf{Role} & \textbf{Example} \\
      \midrule
      Web browser & Firefox, Google Chrome, Microsoft Edge \\
      PDF reader & Adobe Acrobat Reader \\
      Remote-support tool & TeamViewer, AnyDesk, Chrome Remote Desktop \\
      \bottomrule
    \end{tabular}
  \end{center}
  \vspace{0.25cm}
  \muted{Match the software to its job, not to a similar-sounding name.}
\end{frame}

\begin{frame}{Associations at a Glance}
  \begin{center}
    \begin{tabular}{p{0.26\textwidth}p{0.60\textwidth}}
      \toprule
      \textbf{Name} & \textbf{Correct association} \\
      \midrule
      PNG & An image format (Portable Network Graphics) \\
      PDF & A fixed-layout document format by Adobe \\
      \texttt{.mdb} / \texttt{.accdb} & A Microsoft Access database file \\
      Firefox & A web browser by Mozilla \\
      PARAM & A supercomputer by C-DAC \\
      \texttt{.xlsx} & A Microsoft Excel workbook \\
      \bottomrule
    \end{tabular}
  \end{center}
\end{frame}

\begin{frame}{Trap Alert: Confusions to Avoid}
  \begin{itemize}
    \item \texttt{.png} is an \key{extension}; \key{PNG} is the \key{format}
      (TRAP-34).
    \item Firefox is a \key{browser by Mozilla}, not a PDF reader or a
      Microsoft product.
    \item \key{PARAM} is a supercomputer by \key{C-DAC}, not a file format.
    \item \texttt{.mdb} and \texttt{.accdb} are \key{both Access databases}.
    \item Adobe Acrobat Reader reads PDFs; it is \key{not} a web browser.
  \end{itemize}
\end{frame}

\begin{frame}{Lesson 55: Recall}
  \begin{itemize}
    \item An \key{extension} is the letters after the dot; a \key{format} is
      how the data is stored.
    \item PNG = Portable Network Graphics, an image format; JPEG and GIF are
      image formats too.
    \item PDF = Portable Document Format, a fixed-layout format by Adobe;
      scanned PDFs need OCR.
    \item Office family: \texttt{.docx} (Word), \texttt{.xlsx} (Excel),
      \texttt{.pptx} (PowerPoint).
    \item CSV = Comma-Separated Values; TXT = plain text with no formatting.
    \item \texttt{.mdb} and \texttt{.accdb} are both Microsoft Access
      databases.
    \item Firefox is a browser by Mozilla; PARAM is a supercomputer by
      C-DAC.
  \end{itemize}
\end{frame}
\end{document}
